\chapter{Torsional junction of a trapped region and construction of a descending cosmological branch}
\label{chap:universo-descendiente}

\section{Torsional gluing of a trapped region: exact, conditioned, and physical strata of the descendant cosmological branch}

The \cref{ch:rebote-ec} establishes the memory law
\begin{equation}
 M_n=\left(\frac{a_n}{a_0}\right)^{-6}
 \label{eq:descendant-memory-law}
\end{equation}
and formulates its realization through a quadratic spin density in
Einstein--Cartan. The present chapter studies the additional passage from a
torsional bounce to an expansive causal region situated beyond the trapped
interior of a black hole.

The exposition distinguishes \(3\) strata:
\begin{enumerate}
 \item \emph{structure exact}: orientation TRIT, involution of Möbius,
       conservation isometric of the canal and law \(V^{-2}\);
 \item \emph{conditioned consequence}: Einstein--Cartan bounce,
       geometric splice and causal opening under declared hypotheses;
 \item \emph{classified geometric realization}: explicit existence of
       the descendant continuation, minimal, and obstruction to the classical
       branched cobordism with topological change.
\end{enumerate}
This stratification makes it possible to develop the complete image and preserve the
status of each assertion.

The common antecedent is an enriched state
\begin{equation}
 X^{\rm enr}_{\TPK}
 =(x,\text{fiber},\text{orientation},\text{carry},
   \text{boundary},\text{route},\text{memory},\text{observer}),
 \label{eq:descendant-enriched-state}
\end{equation}
produced by \(\mathrm{APP}\to\TRIT\to\TPK\). The causal,
torsional, topological, and quantum readers are applied to this same state and
preserve its genealogy.

\section{Causal orientation of the \(3\) leaves}

The TRIT grammar is realized through the algebras
\begin{equation}
 \mathbb A_\tau=\R[X]/(X^2+\tau),
 \qquad
 J_\tau^2=-\tau I,
 \qquad
 \tau\in\{+1,0,-1\}.
 \label{eq:descendant-TRIT-algebras}
\end{equation}
Their exponentials are
\begin{align}
 \ee^{tJ_{+1}}&=\cos t\,I+\sin t\,J_{+1},
 \label{eq:descendant-elliptic}\\
 \ee^{tJ_0}&=I+tJ_0,
 \label{eq:descendant-parabolic}\\
 \ee^{tJ_{-1}}&=\cosh t\,I+\sinh t\,J_{-1}.
 \label{eq:descendant-hyperbolic}
\end{align}
The sheet \(+1\) carries closure, return, and memory; the sheet \(0\) carries the threshold,
present, and first variation; the sheet \(-1\) carries openness, propagation, and
the two-way future.

Let \(\vartheta\) be an oriented expansion scalar. Define the
causal reader
\begin{equation}
 \boxed{
 \tau_{\rm cau}(\vartheta):=-\operatorname{sgn}(\vartheta)
 \in\{-1,0,+1\}.}
 \label{eq:descendant-causal-reader}
\end{equation}
A positive expansion publishes opening \(-1\); a zero expansion
publishes threshold \(0\); a negative expansion publishes closure \(+1\).

In to family of spheres, \(\vartheta=\theta_{(\ell)}\) may tomarse as
the expansion null saliente. In a region homogeneous or anisotropic,
\(\vartheta=\Theta\) may tomarse as the expansion volumetric. These
two observables occupy distinct domains and share the same
sign rule.

\begin{proposition}[Double-threshold causal sequence]
\label{prop:descendant-double-threshold}
Assume the sign conditions
\begin{equation}
 \theta_{(\ell)}>0\ \text{in the exterior},
 \quad
 \theta_{(\ell)}=0\ \text{at the horizon},
 \quad
 \theta_{(\ell)}<0\ \text{in the trapped interior},
 \label{eq:descendant-null-signs}
\end{equation}
and
\begin{equation}
 \Theta<0\ \text{before the bounce},
 \quad
 \Theta=0\ \text{at the bounce},
 \quad
 \Theta>0\ \text{after the bounce}.
 \label{eq:descendant-volume-signs}
\end{equation}
The reader \cref{eq:descendant-causal-reader} produces the sequence
\begin{equation}
 \boxed{
 -1_{\rm exterior}
 \longrightarrow0_{\rm horizonte}
 \longrightarrow+1_{\rm interior}
 \longrightarrow0_{\rm rebote}
 \longrightarrow-1_{\rm descendiente}.}
 \label{eq:descendant-five-stage-sequence}
\end{equation}
\end{proposition}

\begin{proof}
The assertion follows by applying
\(\tau_{\rm cau}(\vartheta)=-\operatorname{sgn}(\vartheta)\) to each sign of
\cref{eq:descendant-null-signs,eq:descendant-volume-signs}.
\end{proof}

The two states \(0\) are realizations of the same threshold archetype. The
first expresses null marginality; the second expresses volumetric marginality.
The identity of the sign preserves the grammar, while the
declared observable determines the physical content of each threshold.

\section{Anisotropic interior and quadratic memory}

An interior region, spherically symmetric, admits the Kantowski--Sachs
reduction
\begin{equation}
 \dd s^2=-N(\tau)^2\dd\tau^2
 +a_\parallel(\tau)^2\dd\chi^2
 +a_\perp(\tau)^2\dd\Omega_2^2.
 \label{eq:descendant-KS-metric}
\end{equation}
The volume of a comoving cell is
\begin{equation}
 V(\tau)=V_0a_\parallel a_\perp^2,
 \label{eq:descendant-KS-volume}
\end{equation}
and its expansion is
\begin{equation}
 \Theta
 :=\frac1N\frac{\dot V}{V}
 =H_\parallel+2H_\perp,
 \qquad
 H_i:=\frac1N\frac{\dot a_i}{a_i}.
 \label{eq:descendant-KS-expansion}
\end{equation}

A comoving spin charge \(N_s\) produces a density
\begin{equation}
 n_s=\frac{N_s}{V},
 \label{eq:descendant-spin-number-density}
\end{equation}
and the Cartan equation, linear in the current, generates a correction of quadratic form
\begin{equation}
 \rho_s=s_0\left(\frac V{V_0}\right)^{-2}.
 \label{eq:descendant-spin-density}
\end{equation}
The Perron branch publishes exactly the same covariance:
\begin{equation}
 M_n=V_n^{-2}.
 \label{eq:descendant-Perron-volume}
\end{equation}
The torsional transducer
\begin{equation}
 \mathcal T_{\rm EC}:M_n\longmapsto
 \rho_{s,n}=s_0M_n
 \label{eq:descendant-EC-transducer}
\end{equation}
preserves positivity and the volume exponent. The amplitude \(s_0\) follows
from the Cartan current and retains its obligation of prospective selection.

For the discussion local adoptamos units geometric \(c=1\). A
restriction efectiva of Kantowski--Sachs has the form
\begin{equation}
 \frac{\Theta^2}{3}
 =8\pi G(\rho_m-\rho_s)
 +\sigma^2-\frac12\,{}^{(3)}R+\Lambda,
 \label{eq:descendant-KS-constraint}
\end{equation}
where \(\sigma^2\) is the shear and \({}^{(3)}R\) the curvature of the sheet.
A regular volumetric bounce satisfies
\begin{equation}
 \boxed{
 \Theta(\tau_b)=0,
 \qquad
 \frac1N\dot\Theta(\tau_b)>0,}
 \label{eq:descendant-bounce-conditions}
\end{equation}
together with positive finiteness of \(a_\parallel,a_\perp\) and of the curvature
and torsion invariants.

\begin{criterion}[Anisotropic torsional dominance]
\label{crit:descendant-torsion-dominance}
Suppose that, during the contraction,
\begin{equation}
 \rho_s=s_0V^{-2},
 \qquad
 \sigma^2=\Sigma_0^2V^{-2}+o(V^{-2}),
 \label{eq:descendant-torsion-shear-scaling}
\end{equation}
and that matter, curvature, and \(\Lambda\) grow with a power strictly
less than \(V^{-2}\). The inequality
\begin{equation}
 \boxed{8\pi Gs_0>\Sigma_0^2}
 \label{eq:descendant-torsion-dominates-shear}
\end{equation}
is a sufficient condition for the torsional term to dominate the
shear in the regime of small volume. The existence of a transverse root
of \cref{eq:descendant-KS-constraint} and the inequality
\cref{eq:descendant-bounce-conditions} complete the bounce criterion.
\end{criterion}

Asymptotic dominance supplies the mechanism; the transversal root supplies
the geometric event. This distinction preserves the dependence on
curvature, anisotropy, equation of state, and boundary data.

\section{Nonadic condition for the bounce level}

Write \(x:=a/a_0\) for the dimensionless isotropic scale. At the
complete returns \(n=9q\), the Perron branch satisfies
\begin{equation}
 x_{9q}=\frac{a_{9q}}{a_0}=\varphi_{\HMT}^{3q}.
 \label{eq:descendant-nonadic-scale}
\end{equation}
In the isotropic reduction with
\begin{equation}
 \rho_m=\rho_0x^{-3(1+w)},
 \qquad
 \rho_s=s_0x^{-6},
 \qquad
 w<1,
 \label{eq:descendant-isotropic-densities}
\end{equation}
the equality of densities fixes
\begin{equation}
 x_b^{3(1-w)}=\frac{s_0}{\rho_0}.
 \label{eq:descendant-isotropic-bounce}
\end{equation}
If the bounce is required over a complete TPK turn, the elimination of
\(x_b\) yields
\begin{equation}
 \boxed{
 q_b=
 \frac{\log(s_0/\rho_0)}
 {9(1-w)\log\varphi_{\HMT}},
 \qquad
 q_b\in\Z.}
 \label{eq:descendant-quantized-bounce}
\end{equation}
The integrality of \(q_b\) is an additional condition of synchronization,
derived from the hypothesis that the bounce coincides with the nonadic return.
A continuous bounce between two circuits satisfies
\cref{eq:descendant-isotropic-bounce} and has a different status.

\begin{proposition}[Falsifiability of the Synchronous Bounce]
\label{prop:descendant-quantization-falsifier}
With \((s_0,\rho_0,w)\) fixed prospectively, the synchronized realization
exists only if the number of
\cref{eq:descendant-quantized-bounce} is an integer. The distance
\begin{equation}
 \delta_9:=\operatorname{dist}(q_b,\Z)
 \label{eq:descendant-quantization-defect}
\end{equation}
is a quantitative falsifier and cancels exactly at an allowed level.
\end{proposition}

\begin{proof}
The condition \(n=9q\) requires \(q\in\Z\). Substituting
\cref{eq:descendant-nonadic-scale} into
\cref{eq:descendant-isotropic-bounce} and taking logarithms produces the formula.
The characterization of \(\delta_9\) follows from the definition of distance to
the integers.
\end{proof}

\section{Matching data on the bounce hypersurface}

Let \(\Sigma_b\) be a spatial hypersurface separating a contractive region \(\mathcal M_-\) from an expansive region \(\mathcal M_+\). Denote by \(h_{ij}\) the induced metric, by \(n^\mu\) the oriented normal, and by \(\mathring K_{ij}^\pm\) the extrinsic curvatures computed with the Levi--Civita connection on each side. We adopt
\begin{equation}
 [X]:=X^+\big|_{\Sigma_b}-X^-\big|_{\Sigma_b},
 \qquad
 \kappa:=\frac{8\pi G}{c^4}.
 \label{eq:descendant-jump-kappa}
\end{equation}

The Israel metric conditions are
\begin{align}
 [h_{ij}]&=0,
 \label{eq:descendant-Israel-first}\\
 [\mathring K_{ij}-h_{ij}\mathring K]
 &=-\kappa S^{\Sigma}_{ij},
 \label{eq:descendant-Israel-second}
\end{align}
where \(S^\Sigma_{ij}\) is the surface energy--momentum tensor under the fixed-normal
convention. A junction without a material layer satisfies
\begin{equation}
 [h_{ij}]=0,
 \qquad
 [\mathring K_{ij}]=0.
 \label{eq:descendant-shell-free-Israel}
\end{equation}
The inversion of \(\Theta\) then crosses to continuous zero and remains
realizada by the dynamics, rather than being introduced by to jump of the
normal.

In the first-order formulation, let
\(\omega^{ab}=\mathring\omega^{ab}+C^{ab}\) be the connection with contorsion,
and write \(T_\mu:=T^\lambda{}_{\mu\lambda}\) for the trace of the
torsion.
We fix the normalization of the spin current through the Cartan equation
\begin{equation}
 \mathscr C(T)^{\lambda}{}_{\mu\nu}
 :=T^{\lambda}{}_{\mu\nu}
 +\delta^\lambda_\mu T_\nu
 -\delta^\lambda_\nu T_\mu
 =\kappa s^{\lambda}{}_{\mu\nu}.
 \label{eq:descendant-Cartan-equation}
\end{equation}
Its distributional part over \(\Sigma_b\) requires
\begin{equation}
 \boxed{
 \mathscr C(T_\Sigma)
 =\kappa s_\Sigma.}
 \label{eq:descendant-Cartan-junction}
\end{equation}
In the absence of a superficial spin layer,
\(s_\Sigma=0\), the admissible connection lacks a delta torsional component,
and the tangential contorsion has regular gluing. An action with a
Holst term or nonminimal fermionic couplings replaces
\(\mathscr C\) by the corresponding algebraic operator; the conditions
are obtained by integrating that equation across the layer.

\begin{criterion}[Regular Israel--Cartan gluing]
\label{crit:descendant-Israel-Cartan}
A regular torsional seam through \(\Sigma_b\) requires:
\begin{enumerate}[label=\roman*)]
 \item continuity of the induced costructure, equivalent to
       \cref{eq:descendant-Israel-first} up to Lorentz gauge;
 \item equation of Israel \cref{eq:descendant-Israel-second} with every capa
       superficial declared;
 \item Cartan equation \cref{eq:descendant-Cartan-junction} with every
       declared surface current;
 \item continuity of the gauge charges and of the route orientation;
 \item finiteness of the curvature and torsion invariants.
\end{enumerate}
\end{criterion}

The Ashtekar--Barbero Connection
\begin{equation}
 A_i^a=\Gamma_i^a+\gamma_{\rm BI}K_i^a
 \label{eq:descendant-Barbero-connection}
\end{equation}
can act as a reader of the seam when both sides share
costructure, \(\gamma_{\rm BI}\), and normalization. Its continuity is a
conditional consequence of the corresponding gluing data; its
verification requires the explicit transport of \(\Gamma\) and \(K\) to the
same boundary section.

\begin{proposition}[Joint involution of Barbero and extrinsic curvature]
\label{prop:descendant-Barbero-involution}
Consider the algebraic involution
\begin{equation}
 \mathfrak J_b:
 (\Gamma_i^a,K_i^a,\gamma_{\rm BI})
 \longmapsto
 (\Gamma_i^a,-K_i^a,-\gamma_{\rm BI}).
 \label{eq:descendant-Barbero-involution}
\end{equation}
Then
\begin{equation}
 \mathfrak J_b(A_i^a)=A_i^a,
 \label{eq:descendant-Barbero-A-invariant}
\end{equation}
and every area spectrum that depends on \(|\gamma_{\rm BI}|\) remains
invariant. Identification of \(\mathfrak J_b\) with the dynamic inversion
of the bounce requires the Israel--Cartan intertwiner to transport
the orientation of \(K\) and the sheet of \(\gamma_{\rm BI}\) according to
\cref{eq:descendant-Barbero-involution}.
\end{proposition}

\begin{proof}
The product of the two oriented factors satisfies
\((-\gamma_{\rm BI})(-K_i^a)=\gamma_{\rm BI}K_i^a\), while \(\Gamma_i^a\)
remains fixed. The areal dependence on \(|\gamma_{\rm BI}|\) is even.
\end{proof}

\section{Null horizon as a marginal hypersurface and initial matching condition}

The horizon hypersurface \(\Sigma_H\) has null character and requires
adapted junction data. Let \(q_{AB}\) be the two-dimensional metric of the
sections, \(\ell^\mu\) the null generator, and \(N^\mu\) a transverse vector with
\(\ell\cdot N=-1\). The transverse curvature is
\begin{equation}
 \mathcal C_{AB}
 :=q_A{}^\mu q_B{}^\nu\nabla_\mu N_\nu.
 \label{eq:descendant-null-transverse-curvature}
\end{equation}
A null junction without an impulsive layer preserves \(q_{AB}\), the class of
\(\mathcal C_{AB}\) that determines the surface tension and the tangential
Cartan flux. Marginality
\begin{equation}
 \theta_{(\ell)}\big|_{\Sigma_H}=0
 \label{eq:descendant-null-marginality}
\end{equation}
realizes the first state \(0\) of
\cref{eq:descendant-five-stage-sequence}. The formulation null mantiene
separada this condition of the condition volumetric
\(\Theta|_{\Sigma_b}=0\).

\section{Möbius band and orientation memory}

The dodecaphasic words are acted on by the involution of genealogical type
\begin{equation}
 C_M(\mathbf t):=-\operatorname{rev}(\mathbf t),
 \qquad
 C_M^2=I.
 \label{eq:descendant-CM}
\end{equation}
The associated oriented line is obtained by the quotient
\begin{equation}
 (\mathbf t,\chi)\sim(C_M\mathbf t,-\chi).
 \label{eq:descendant-Mobius-quotient}
\end{equation}
When the base is a return loop, this identification defines a band
with holonomy \(-I\): one turn reverses orientation and two turns
restore it.

Let \(\mathscr H_\gamma\) denote the holonomy over one turn. The hierarchy
\begin{equation}
 \mathcal R_{\rm vec}
 \subseteq\mathcal R_{\rm ray}
 \subseteq\mathcal R_{\rm obs}
 \label{eq:descendant-return-hierarchy}
\end{equation}
distinguishes return of vector, return of ray, and return of observable. For
\(\mathscr H_\gamma=-I\), the ray returns in one turn, the oriented vector
returns in two and the sign memory preserves the traversal.

\begin{proposition}[Topological status of the band]
\label{prop:descendant-Mobius-status}
The relation \cref{eq:descendant-Mobius-quotient} defines a real line
that is nonorientable over the loop exactly when its first
Stiefel--Whitney class satisfies
\begin{equation}
 w_1=1\in H^1(S^1;\Z_2).
 \label{eq:descendant-w1}
\end{equation}
This assertion concerns to the haz genealogical of orientation. A
realization as topology of the spacetime requires to morphism toward
the frame or spinor bundle that preserves causality and Cartan structure.
\end{proposition}

\begin{proof}
A real line over \(S^1\) is determined by the parity of its
transition function. The transition \(-1\) of
\cref{eq:descendant-Mobius-quotient} represents the nontrivial class; two
circuits square the transition and yield \(+1\).
\end{proof}

The band provides an exact representation of orientation and memory.
Physical continuity across the horizon and the bounce is formulated
through the junctions of the preceding sections; both structures are
coordinated by an intertwiner and preserve their domains.

\section{Relative Observer, Parallel Transport and Autoinertia in the Descending Cosmological Region}

The expanded dynamic state includes observer and memory:
\begin{equation}
 \Gamma=(x,\mathcal O,m),
 \label{eq:descendant-observer-state}
\end{equation}
and the frame accessible to the observer is a restriction of the total connection,
\begin{equation}
 \mathfrak F_{\mathcal O}(\Gamma)
 :=\operatorname{Res}_{\mathcal O}
 (\Gamma,\nabla_{\TPK}).
 \label{eq:descendant-observer-frame}
\end{equation}
A horizontal trajectory of the complete system,
\begin{equation}
 \nabla^{\rm tot}_{\dot\Gamma}\dot\Gamma=0,
 \label{eq:descendant-total-horizontal}
\end{equation}
can publish acceleration in a subsystem \(S\):
\begin{equation}
 a_S=\nabla^S_{\dot\gamma_S}\dot\gamma_S
 =\mathcal K_{\pi_S}(\dot\Gamma,\dot\Gamma),
 \label{eq:descendant-projected-acceleration}
\end{equation}
where \(\mathcal K_{\pi_S}\) is the second fundamental tensor of the
projection. Self-inertia denotes the horizontality of the whole and the
response induced in its projections.

The coupling between system, apparatus, and register is described by an
isometry
\begin{equation}
 V_{\rm obs}:\mathcal H_{\rm in}\longrightarrow
 \mathcal H_{\rm sys}\otimes\mathcal H_{\rm mem},
 \qquad
 V_{\rm obs}^*V_{\rm obs}=I.
 \label{eq:descendant-observer-isometry}
\end{equation}
Let \(\{P_y\}\) be the PVM of the register in \(\mathcal H_{\rm mem}\), with
\begin{equation}
 P_yP_z=\delta_{yz}P_y,
 \qquad
 \sum_yP_y=I_{\rm mem}.
 \label{eq:descendant-observer-PVM}
\end{equation}
The isometry and this PVM induce the completely positive instrument
\begin{equation}
 \mathcal I_y(\rho)
 :=\Tr_{\rm mem}\!\left[
 (I_{\rm sys}\otimes P_y)V_{\rm obs}\rho V_{\rm obs}^*
 (I_{\rm sys}\otimes P_y)
 \right],
 \qquad
 \sum_y\Tr\mathcal I_y(\rho)=\Tr\rho.
 \label{eq:descendant-observer-instrument}
\end{equation}
The relative states are then obtained by means of
\begin{equation}
 p(y)=\Tr\mathcal I_y(\rho),
 \qquad
 \rho_y=\frac{\mathcal I_y(\rho)}{p(y)}.
 \label{eq:descendant-relative-state}
\end{equation}
The section \(\rho_y\) is the state relative to the record \(y\); the global state
preserves all channels in the image of the isometry. This
structure formalizes the kinship with Everett at the precise level of
relative states and leaves the ontology of the branches as a choice in subsequent physical terms.

\section{Global channel of the descendant region}

The operatorial version of a descendant opening uses an isometry
\begin{equation}
 V_{\rm pc}:\mathcal H_{\rm in}\longrightarrow
 \mathcal H_{\rm parent}\otimes\mathcal H_{\rm child},
 \qquad
 V_{\rm pc}^*V_{\rm pc}=I.
 \label{eq:descendant-parent-child-isometry}
\end{equation}
The joint channel and its complementary channels are
\begin{align}
 \mathcal N_{\rm pc}(\rho)
 &=V_{\rm pc}\rho V_{\rm pc}^*,
 \label{eq:descendant-global-channel}\\
 \mathcal N_{\rm parent}(\rho)
 &=\Tr_{\rm child}\mathcal N_{\rm pc}(\rho),
 \qquad
 \mathcal N_{\rm child}(\rho)
 =\Tr_{\rm parent}\mathcal N_{\rm pc}(\rho).
 \label{eq:descendant-complementary-channels}
\end{align}

\begin{theorem}[Global Conservation and Relative Distribution]
\label{thm:descendant-global-conservation}
The map \(\mathcal N_{\rm pc}\) preserves trace and purity of to pure state
on its image. The parental and descendant channels may produces
mixed states, and for to pure bipartite output satisfy
\begin{equation}
 S(\rho_{\rm parent})=S(\rho_{\rm child}).
 \label{eq:descendant-equal-entanglement}
\end{equation}
The global information is contained in the outputs and their correlations.
\end{theorem}

\begin{proof}
The identity \(V_{\rm pc}^*V_{\rm pc}=I\) implies
\[
 \Tr(V_{\rm pc}\rho V_{\rm pc}^*)
 =\Tr(\rho V_{\rm pc}^*V_{\rm pc})=\Tr\rho.
\]
An isometry maps unit vectors to unit vectors. Equality of the
reduced entropies follows from the Schmidt decomposition of the pure
bipartite state.
\end{proof}

The conservation of a charge \(Q\) or of energy requires, in addition,
the intertwinings
\begin{equation}
 Q_{\rm out}V_{\rm pc}=V_{\rm pc}Q_{\rm in},
 \qquad
 H_{\rm out}V_{\rm pc}=V_{\rm pc}H_{\rm in}.
 \label{eq:descendant-charge-energy-intertwining}
\end{equation}
The isometry preserves probability; the equations
\cref{eq:descendant-charge-energy-intertwining} preserve the declared dynamic
quantities. An admissible output distributes the information among
correlations and respects the impossibility of cloning an arbitrary state into
two identical copies.

The provenance of the sectors \(90/120\) is preserved through the modular
holographic Hamiltonian
\begin{equation}
 K_{\rm HHM}=c_0I+\Delta_{\rm mod}
 (90N_{90}+120N_{120}).
 \label{eq:descendant-HHM}
\end{equation}
The total occupation and the weighted memory,
\begin{equation}
 N=N_{90}+N_{120},
 \qquad
 D=90N_{90}+120N_{120},
 \label{eq:descendant-HHM-data}
\end{equation}
Transport of the total occupation is imposed as an independent operator identity:
\begin{equation}
 N^{\rm out}V_{\rm pc}=V_{\rm pc}N^{\rm in}.
 \label{eq:descendant-number-intertwining}
\end{equation}
With \(N\) and the weighted memory \(D\) conserved, the two sectors are
reconstructed uniquely:
\begin{equation}
 N_{90}=4N-\frac D{30},
 \qquad
 N_{120}=\frac D{30}-3N.
 \label{eq:descendant-HHM-recovery}
\end{equation}
The condition of transport
\begin{equation}
 K_{\rm HHM}^{\rm out}V_{\rm pc}
 =V_{\rm pc}K_{\rm HHM}^{\rm in}
 \label{eq:descendant-HHM-intertwining}
\end{equation}
The identities
\cref{eq:descendant-number-intertwining,eq:descendant-HHM-intertwining}
jointly preserve occupation and sectorial provenance across the
seam. These equalities are added to the dynamic obligations of a
complete realization.

\section{Descendant continuation and branched cobordism}

The denomination \emph{descendant universe} admits two
geometric levels.

\begin{definition}[Minimal descendant continuation]
\label{def:descendant-minimal}
A descendant continuation is a region
\(\mathcal M_{\rm child}\) situated to the future of the rebound, with
\begin{equation}
 \Theta>0,
 \qquad
 \text{finite curvature and torsion},
 \qquad
 \text{a global time function in its causal domain}.
 \label{eq:descendant-minimal-conditions}
\end{equation}
The region may pertenecer to a prolongation conectada of the interior and
preserves the topology of the hojas espaciales.
\end{definition}

\begin{definition}[Parent-Child Cobordism]
\label{def:descendant-cobordism}
The branched version is an effective cobordism
\begin{equation}
 \mathcal W:\Sigma_{\rm in}\longrightarrow
 \Sigma_{\rm parent}\sqcup\Sigma_{\rm child},
 \label{eq:descendant-cobordism}
\end{equation}
equipped with Cartan structure, gluing data
\cref{crit:descendant-Israel-Cartan}, and channel
\cref{eq:descendant-parent-child-isometry}.
\end{definition}

The minimal continuation requires a bounce and a causal extension. The
branched cobordism adds topological change and separation of outputs.
In a Lorentzian geometry that is classical, nondegenerate, and globally
hyperbolic, topological change is subject to causal obstructions.
The branched version must therefore be realized through a region that is
quantum or torsional and in which it is declared which classical hypothesis is modified,
or through an interpretation of two causal domains within a single
global topology.

\begin{criterion}[Promotion to a Physical Descendant Universe]
\label{crit:descendant-promotion}
The complete image acquires the status of a realization that is physical when the following are
constructed simultaneously:
\begin{enumerate}[label=\roman*)]
 \item a solution of the Einstein--Cartan equations with
       \cref{eq:descendant-bounce-conditions};
 \item regular Israel--Cartan data on \(\Sigma_H\) and \(\Sigma_b\);
 \item a future expansive region with temporal function and Cauchy surfaces in
       each accessible branch;
 \item finiteness of invariants and prolongation of geodesics through the
       nucleus;
 \item absence of closed causal curves and unregistered layers;
 \item an isometric channel that preserves the declared charges and the
       Hamiltonian of provenance \cref{eq:descendant-HHM-intertwining};
 \item nonadic naturality of the torsional current and the
       boundary data;
 \item an intertwiner between the genealogical band and the physical bundle of
       frames or spinors.
\end{enumerate}
\end{criterion}

These conditions transforman the narrative
\cref{eq:descendant-five-stage-sequence} in a problema of existence
geometric, operatorial and causal with controles separate.

\section{Local consequences of the torsional matching: conservation, regularity, and causal prolongation}

\begin{proposition}[Content of a Complete Realization]
\label{prop:descendant-conditioned-consequences}
Suppose that the
\cref{crit:descendant-promotion} is satisfied. Then:
\begin{enumerate}[label=\alph*)]
 \item the causal sign traverses the sequence
       \cref{eq:descendant-five-stage-sequence};
 \item the torsional memory preserves the covariance \(V^{-2}\) on both
       sides of the rebound;
 \item the global state evolves through a trace-preserving channel;
 \item each observer accesses a relative state of its branch;
 \item the charges included in
       \cref{eq:descendant-charge-energy-intertwining} are conserved;
 \item the orientation of Möbius is transported as holonomy data.
\end{enumerate}
\end{proposition}

\begin{proof}
Each assertion is the application of one of the conditions of the criterion:
the causal reader proves (a); the torsional transducer proves (b); the isometry
proves (c) and (d); the intertwinings prove (e); and the sheaf morphism proves
(f).
\end{proof}

The proposition organizes consequences of a constructed structure. The
existence of the minimal continuation is resolved by the
\cref{thm:ivv-descendant-existence}; classical branching with a
topological change is classified by the obstruction of the
\cref{thm:ivv-cobordism-obstruction}.

\section{Torsional dynamics of bounce: exact identities, cosmological interfaces and closure conditions}

\begin{center}
\small
\begin{tabularx}{\textwidth}{>{\bfseries\raggedright\arraybackslash}p{31mm}YY}
\toprule
Object & Status & Obligation or control\\
\midrule
Law \(M_n=V_n^{-2}\) & Exact internal result & Perron's Certificate and Refinement\\
Reader \(-\operatorname{sgn}\vartheta\) & Exact formalization & Observable of expansion tipado\\
Homogeneous bounce & Consecuencia condicionada & \(s_0>0\), \(w<1\), transverse zero\\
Kantowski--Sachs Bounce & Physical Interface & Dominance of controlled torsion, shear, and curvature\\
Israel--Cartan Joining & Physical Interface & Declared layer and spin current data\\
Möbius band & Exact Topological Result & Intertwiner toward the physical bundle\\
Estado relativo & Exact operational result & Instrument and observed fiber\\
Parent--offspring canal & Conditional formalization & Isometry and Conservation of Charges\\
Descendant continuation & Existence theorem & Solution regular and causally complete\\
Cobordismo ramificado & Proven classical obstruction & Incompatibility with global hyperbolicity and topological change\\
\bottomrule
\end{tabularx}
\end{center}

\section{Obstructions to Cosmological Continuation: Loss of Orientation, Junction Regularity, or Return Memory}

The negative controls distinguish each layer:
\begin{enumerate}
 \item a solution with \(\Theta=0\) and \(\dot\Theta<0\) represents a
       volume maximum and remains outside the bounce;
 \item the case \(8\pi Gs_0\leq\Sigma_0^2\) removes the guarantee of
       anisotropic torsional dominance;
 \item a \(q_b\notin\Z\) preserves the continuous rebound and refutes the version
       synchronized with a complete turn;
 \item a jump of \(h_{ij}\), of the Israel quantity, or of the
       Cartan current requires an explicit surface layer;
 \item a band with \(w_1=0\) has global orientation and refutes the
       realization of Möbius;
 \item a map \(V_{\rm pc}\) with \(V_{\rm pc}^*V_{\rm pc}\ne I\) loses
       global probability;
 \item two identical copies of an arbitrary state violate the no-
       cloning control and refute the correlated distribution reading;
 \item a descendant region with closed causal curves, divergent invariants,
       or inextendible geodesics fails the causal promotion;
 \item a topological change that simultaneously preserves all the
       obstructed classical hypotheses lacks the required geometric
       domain.
\end{enumerate}

\begin{proofstatusbox}
The grammar \(-1,0,+1\), the sign reader, the involution \(C_M^2=I\), the classification of the band, the law \(V^{-2}\), the hierarchy of returns, and conservation of an isometric channel are exact results in their domains. The Einstein--Cartan bounce is a consequence of the declared equations and inequalities. The extension to Kantowski--Sachs, the junctions, continuity of the Barbero connection, and causal opening are physical interfaces. The minimal parent--descendant continuation is realized by the \cref{thm:ivv-descendant-existence}; branched topological change under the classical hypotheses is excluded by the \cref{thm:ivv-cobordism-obstruction}.
\end{proofstatusbox}

\begin{falsifierbox}
The torsional seam is refuted if the current constructed from HMT produces a power different from \(V^{-2}\), an amplitude with an incompatible sign, or a failure of refinement. The bounce is refuted by the absence of a regular transverse root. The junction is refuted by an Israel or Cartan residue without a surface source. Global conservation is refuted by loss of trace, norm, or an intertwined charge. Promotion to a descendant universe is refuted by a singularity, causal violation, or the impossibility of constructing the required orientation bundle and global channel. Each refutation affects its interface and preserves the preceding internal results.
\end{falsifierbox}
